\documentclass[11pt]{amsart}

\usepackage[margin=1in]{geometry}
\usepackage[T1]{fontenc}
\usepackage{lmodern}
\usepackage{amsmath,amssymb,amsfonts,amsthm,mathtools}
\usepackage{enumitem}
\usepackage[colorlinks=true,linkcolor=blue,citecolor=blue,urlcolor=blue]{hyperref}
\usepackage[noabbrev,capitalize]{cleveref}
\setlength{\emergencystretch}{2em}

\newtheorem{theorem}{Theorem}[section]
\newtheorem{proposition}[theorem]{Proposition}
\newtheorem{lemma}[theorem]{Lemma}
\newtheorem{corollary}[theorem]{Corollary}
\theoremstyle{definition}
\newtheorem{definition}[theorem]{Definition}
\newtheorem{assumption}[theorem]{Assumption}
\theoremstyle{remark}
\newtheorem{remark}[theorem]{Remark}

\newcommand{\R}{\mathbb R}
\newcommand{\Mcal}{\mathcal M}
\newcommand{\Etheta}{\mathcal E^\theta}
\newcommand{\Hcal}{\mathcal H}
\newcommand{\Lcal}{\mathcal L}
\newcommand{\Gcal}{\mathcal G}
\newcommand{\Dcal}{\mathcal D}
\newcommand{\Tr}{\operatorname{tr}}

\title[Representation Calculus for theta-Expectations]{Representation Calculus for theta-Expectations}
\author{Qian Qi}
\date{\today}

\begin{document}

\begin{abstract}
This paper develops the representation theory that lies downstream of the
first-principles theta-expectation construction.  The nonlinear semigroup and
its nonconvex HJB generator are taken as outputs of the companion
homogenization paper.  After this input is fixed, each sufficiently regular
payoff determines a linearized parabolic generator, a carre-du-champ, a
calibrated diffusion representation, a decoupled BSDE, and a semigroup-level
Cameron-Martin--Girsanov formula.  These objects represent the already-derived
semigroup; they are not inputs to the deterministic billiard derivation.
\end{abstract}

\maketitle

\section{Introduction}

This is the third article in the sequence.  Paper 1 proves the singular
billiard response package.  Paper 2 uses that package to derive the nonconvex
HJB equation and the theta-expectation semigroup from deterministic billiards.
The present paper begins only after those two steps are complete.

\paragraph{Representation, not derivation.}
The central point is the payoff dependence of every classical probabilistic
representation.  A genuinely nonlinear HJB semigroup cannot be represented by a
single classical martingale problem for all test functions.  Once a terminal
payoff is fixed, however, its decoupling field freezes the nonlinear generator
into an ordinary linear parabolic operator.  Diffusions, BSDEs, and Girsanov
formulae are then available for that calibrated linear problem.

\subsection{Theorem map}

\begin{center}
\small
\begin{tabular}{p{0.24\linewidth}p{0.38\linewidth}p{0.28\linewidth}}
\hline
Layer & Paper theorem & Proof role \\
\hline
Semigroup input & \cref{ass:theta_semigroup_input} &
imported from Paper 2 \\
Obstruction & \cref{thm:linearity_obstruction} &
rules out one universal law \\
Linearization & \cref{thm:linearized_representation} &
freezes a fixed payoff \\
Regularized diffusion & \cref{thm:regularized_diffusion_representation} &
represents the calibrated linear problem \\
BSDE & \cref{thm:fbsde_representation} &
uses \(Z=\sigma^\top\nabla u\) \\
Orientation & \cref{prop:parabolic_orientation_ledger} &
checks terminal-value signs \\
Girsanov & \cref{thm:nonlinear_cm_formula} &
semigroup-level shift identity \\
Stress test & \cref{prop:robust_pricing_stress_test} &
short application illustration \\
\hline
\end{tabular}
\end{center}

\subsection{Provenance and dependency guardrail}

The stochastic-analysis tools used below are classical: the martingale-problem
representation follows \cite{stroock2007multidimensional}, the parabolic and
viscosity stability language follows
\cite{CrandallIshiiLions1992,FlemingSoner2006}, the BSDE convention follows
\cite{PardouxPeng1992}, and the stochastic-calculus/Girsanov notation follows
\cite{protter2012stochastic}.  Nonlinear expectation terminology follows
\cite{Peng2019}.  These references are used only after the deterministic
theta-expectation semigroup has been derived in the companion HJB paper.

\section{Input from the theta-Expectation Paper}

\begin{assumption}[Derived nonlinear semigroup]
\label{ass:theta_semigroup_input}
The companion theta-expectation paper has constructed coefficients
\[
  D(x,p)\in\R^{d\times d}_{\mathrm{sym},+},
  \qquad H(x,p)\in\R,
\]
from deterministic billiard response theory and finite-response mechanics.  It
has also proved comparison and uniqueness for the terminal-value HJB equation
\begin{equation}
\label{eq:input_hjb}
  -\partial_t u-\Tr(D(x,\nabla u)\nabla^2u)-H(x,\nabla u)=0,
  \qquad u(T,\cdot)=\phi .
\end{equation}
The normalized semigroup is
\begin{equation}
  \Etheta_{t,T}[\phi](x)=u(t,x).
\end{equation}
It is constant preserving, cash additive, monotone, and time consistent, and it
is generally nonlinear and not subadditive.
\end{assumption}

\begin{remark}
\Cref{ass:theta_semigroup_input} is the only connection to deterministic
billiards used in this paper.  All microscopic response estimates and the HJB
homogenization theorem belong to the companion papers.
\end{remark}

\begin{assumption}[Representation window]
\label{ass:representation_window}
Fix a compact cylinder \(Q=[t_0,T]\times K\) and a compact jet set containing
all gradients and Hessians used below.  On \(Q\), the payoff under discussion
has a decoupling field \(u\) that is either \(C^{1,2}\) or the locally uniform
limit of smooth uniformly parabolic regularizations.  The coefficients
\(D,H\) and their first \(p\)-derivatives are bounded and continuous on this
jet set.  In the nondegenerate diffusion and BSDE statements, also assume
\(a^u=2D(x,\nabla u)\) is uniformly positive definite and has a chosen square
root \(\sigma^u\) with the regularity needed for the corresponding stochastic
equations.
\end{assumption}

\section{Linearity Obstruction}

\begin{theorem}[Linearity obstruction]
\label{thm:linearity_obstruction}
Let $\Gcal$ be an operator on a vector space of smooth tests.  If there is a
single probability law under which
\begin{equation}
  f(t,X_t)-f(0,X_0)-\int_0^t \Gcal f(s,X_s)\,ds
\end{equation}
is a martingale for every test $f$, then $\Gcal$ is additive on that class.
Consequently a genuinely nonlinear operator
\begin{equation}
  \Gcal f=\partial_t f+\Tr(D(x,\nabla f)\nabla^2f)+H(x,\nabla f)
\end{equation}
cannot be the generator of one classical martingale problem for all tests.
\end{theorem}

\begin{proof}
Apply the martingale property to $f$, $g$, and $f+g$.  Subtracting the three
martingales gives a finite-variation martingale with drift
\[
  \Gcal(f+g)-\Gcal f-\Gcal g .
\]
It is constant, so the drift vanishes along the support of the law.  Varying
the initial point gives additivity on the test class.  Nonlinear dependence on
$\nabla f$ or $\nabla^2 f$ violates this except in accidentally linear cases.
\end{proof}

\section{Payoff-Dependent Linearization}

\begin{definition}[Calibrated linear generator]
\label{def:calibrated_linear_generator}
Work on a representation window satisfying
\cref{ass:representation_window}.  For a fixed terminal payoff \(\phi\), let
\(u\) be its decoupling field.  Set
\begin{align}
  a^u(t,x)&=2D(x,\nabla u(t,x)),\\
  b^u_k(t,x)&=\partial_{p_k}H(x,\nabla u(t,x))
  +\sum_{i,j}\partial_{p_k}D_{ij}(x,\nabla u(t,x))\,\partial_{ij}u(t,x).
\end{align}
The payoff-calibrated generator is
\begin{equation}
  \Lcal_t^u f(x)=\frac12\Tr(a^u(t,x)\nabla^2 f(x))
  +b^u(t,x)\cdot\nabla f(x).
\end{equation}
Its carre-du-champ is
\begin{equation}
  \Gamma_t^u(f,g)
  =\frac12\bigl(\Lcal_t^u(fg)-f\Lcal_t^ug-g\Lcal_t^uf\bigr)
  =\frac12\nabla f^\top a^u\nabla g.
\end{equation}
\end{definition}

\begin{theorem}[Linearized representation for a fixed payoff]
\label{thm:linearized_representation}
Let $u$ be the smooth solution of \eqref{eq:input_hjb} with terminal payoff
$\phi$.  If $u^\delta$ solves the same HJB equation with terminal payoff
$\phi+\delta\eta$ and $u^\delta=u+\delta v+o(\delta)$ locally uniformly, then
$v$ solves
\begin{equation}
  -\partial_t v-\Lcal_t^u v=0,\qquad v(T,\cdot)=\eta .
\end{equation}
Thus the first variation of the nonlinear semigroup at $\phi$ is represented
by the linear parabolic equation generated by $\Lcal^u$.
\end{theorem}

\begin{proof}
Insert $u^\delta=u+\delta v+o(\delta)$ into the HJB operator and differentiate
at $\delta=0$.  The derivative of the Hessian term gives
$\Tr(D(x,\nabla u)\nabla^2 v)$ plus the drift correction
\[
  \sum_{i,j,k}\partial_{p_k}D_{ij}(x,\nabla u)\,\partial_{ij}u\,\partial_k v.
\]
The derivative of $H$ gives $\partial_pH(x,\nabla u)\cdot\nabla v$.  These
terms are exactly $\Lcal_t^u v$ with the sign convention above.
\end{proof}

\section{Calibrated Diffusion Representation}

\begin{theorem}[Smooth nondegenerate calibrated diffusion]
\label{thm:smooth_diffusion_representation}
Assume $a^u$ is uniformly positive definite and the coefficients in
\cref{def:calibrated_linear_generator} are regular enough for the martingale
problem of $\Lcal^u$ to be well posed.  Then the derivative semigroup in
\cref{thm:linearized_representation} has the classical diffusion
representation
\begin{equation}
  dX_s=b^u(s,X_s)\,ds+\sigma^u(s,X_s)\,dW_s,
  \qquad \sigma^u(\sigma^u)^\top=a^u .
\end{equation}
The representing law depends on the payoff through $u$.
\end{theorem}

\begin{proof}
For fixed $u$, the operator $\Lcal^u$ is linear.  The standard parabolic
martingale-problem representation
\cite{stroock2007multidimensional,FlemingSoner2006} applies to that linear
operator.  The
dependence on $\phi$ enters through the coefficients $a^u$ and $b^u$.
\end{proof}

\begin{theorem}[Degenerate representation by viscosity regularization]
\label{thm:regularized_diffusion_representation}
Suppose the viscosity solution $u$ is the locally uniform limit of smooth
solutions $u^\varepsilon$ to uniformly parabolic regularizations with
$D$ replaced by $D+\varepsilon I$.  The corresponding calibrated linear
diffusion semigroups converge locally uniformly to the payoff-calibrated
representation of the original degenerate HJB semigroup.
\end{theorem}

\begin{proof}
For each $\varepsilon>0$, the matrix $2(D+\varepsilon I)$ is uniformly positive
and gives a classical diffusion representation.  Stability of viscosity
solutions \cite{CrandallIshiiLions1992,FlemingSoner2006} and stability of the
linearized coefficients along the smooth regularization identify the locally
uniform limit with the derivative of the original nonlinear semigroup.  The
construction starts after the effective HJB has been derived.
\end{proof}

\section{BSDE Representation}

\begin{lemma}[Gradient variable versus BSDE integrand]
\label{lem:gradient_z_conversion}
In the representation window, the HJB gradient variable is
\[
  p(t,x)=\nabla u(t,x).
\]
The BSDE martingale integrand is not \(p\).  It is
\[
  Z_t=\sigma(X_t,p(t,X_t))^\top p(t,X_t).
\]
When \(\sigma^\top\) is invertible on the chosen nondegenerate jet region, one
may recover \(p=(\sigma^\top)^{-1}Z\).  All BSDE drivers below use this
conversion only after the payoff and decoupling field are fixed.
\end{lemma}

\begin{theorem}[Decoupled FBSDE for a smooth payoff]
\label{thm:fbsde_representation}
Let $u$ be a smooth HJB solution and assume
\[
  \sigma(x,p)\sigma(x,p)^\top=2D(x,p)
\]
is nondegenerate on the relevant jet range.  Let $X$ solve
\begin{equation}
  dX_t=b^u(t,X_t)\,dt+\sigma(X_t,\nabla u(t,X_t))\,dW_t .
\end{equation}
Set
\[
  Y_t=u(t,X_t),\qquad
  Z_t=\sigma(X_t,\nabla u(t,X_t))^\top\nabla u(t,X_t).
\]
Then $(X,Y,Z)$ satisfies the decoupled BSDE
\begin{equation}
  dY_t=-g(t,X_t,Z_t)\,dt+Z_t\,dW_t,
  \qquad Y_T=\phi(X_T),
\end{equation}
where, writing $p=(\sigma^\top)^{-1}Z$,
\begin{equation}
  g(t,x,Z)=H(x,p)-b^u(t,x)\cdot p .
\end{equation}
\end{theorem}

\begin{proof}
Apply Ito's formula to $Y_t=u(t,X_t)$, using the convention in
\cref{lem:gradient_z_conversion}.  The martingale coefficient is
$Z_t=\sigma^\top\nabla u$.  The drift is
\[
  \partial_tu+b^u\cdot\nabla u+\frac12\Tr(\sigma\sigma^\top\nabla^2u)
  =\partial_tu+b^u\cdot p+\Tr(D\nabla^2u).
\]
Using the HJB equation, this equals $b^u\cdot p-H(x,p)=-g(t,x,Z)$.
\end{proof}

\section{Sign Convention and Parabolic Orientation}

\begin{proposition}[Parabolic orientation ledger]
\label{prop:parabolic_orientation_ledger}
The terminal-value convention in \cref{eq:input_hjb}, the calibrated diffusion
orientation, and the BSDE sign convention are consistent as follows.  For
\[
  Y_s=u(s,X_s),\qquad
  Z_s=\sigma(X_s,\nabla u(s,X_s))^\top\nabla u(s,X_s),
\]
the differential form
\[
  dY_s=-g(s,X_s,Z_s)\,ds+Z_s\,dW_s
\]
is equivalent to the backward integral equation
\[
  Y_t=\phi(X_T)+\int_t^T g(s,X_s,Z_s)\,ds
       -\int_t^T Z_s\,dW_s.
\]
If one writes forward parabolic time as \(\tau=T-t\) and
\(U(\tau,x)=u(T-\tau,x)\), then the same terminal-value HJB becomes
\[
  \partial_\tau U
  -\Tr\bigl(D(x,\nabla U)\nabla^2U\bigr)-H(x,\nabla U)=0,
  \qquad U(0,\cdot)=\phi .
\]
Thus no sign is changed when passing from the terminal HJB to the backward
BSDE; only the time orientation changes.
\end{proposition}

\begin{proof}
For the calibrated diffusion, Ito's formula gives
\[
  dY_s
  =
  \bigl(\partial_su+b^u\cdot\nabla u
  +\Tr(D(x,\nabla u)\nabla^2u)\bigr)(s,X_s)\,ds
  +Z_s\,dW_s .
\]
The terminal-value HJB gives
\[
  \partial_su+\Tr(D(x,\nabla u)\nabla^2u)=-H(x,\nabla u),
\]
so the drift is \(b^u\cdot p-H(x,p)=-g(s,x,Z)\).  Integrating the differential
identity from \(t\) to \(T\) gives the displayed BSDE.  The forward-time
equation follows by the substitution \(U(\tau,x)=u(T-\tau,x)\), for which
\(\partial_\tau U=-\partial_tu\).
\end{proof}

\section{Nonlinear Cameron-Martin--Girsanov Formula}

\begin{definition}[Hamiltonian shift]
\label{def:hamiltonian_shift}
For a deterministic shift $h(t,x)$ define
\begin{align}
  H^h(x,p)&=H(x,p+h)-H(x,h),\\
  D^h(x,p)&=D(x,p+h).
\end{align}
The shifted semigroup is generated by the HJB equation with $D,H$ replaced by
$D^h,H^h$.
\end{definition}

\begin{theorem}[Semigroup-level Cameron-Martin formula]
\label{thm:nonlinear_cm_formula}
Let $u_h$ solve the shifted equation.  Then $u_h$ also solves the unshifted
equation with deterministic generator defect
\begin{equation}
\label{eq:lambda_h}
  \Lambda_h(t,x,p,X)
  =\Tr\bigl((D(x,p+h)-D(x,p))X\bigr)
  +H(x,p+h)-H(x,h)-H(x,p).
\end{equation}
Equivalently,
\begin{equation}
  -\partial_tu_h-\Tr(D(x,\nabla u_h)\nabla^2u_h)-H(x,\nabla u_h)
  -\Lambda_h(t,x,\nabla u_h,\nabla^2u_h)=0.
\end{equation}
When $D$ is independent of $p$ and $H$ is linear-quadratic, this identity
reduces after linearization to the classical Cameron-Martin--Girsanov drift
tilt \cite{protter2012stochastic}.
\end{theorem}

\begin{proof}
Insert the definitions of $D^h$ and $H^h$ into the shifted HJB equation.  Add
and subtract the unshifted terms
\[
  \Tr(D(x,\nabla u_h)\nabla^2u_h)+H(x,\nabla u_h).
\]
The remainder is exactly \eqref{eq:lambda_h}.  The linear-quadratic special
case gives
$H(x,p+h)-H(x,h)-H(x,p)=p^\top a(x)h$, the usual drift-tilt term for the
payoff-calibrated linear diffusion.
\end{proof}

\section{A Short Robust-Pricing Illustration}

\begin{definition}[Calibrated stress family]
\label{def:calibrated_stress_family}
Fix a terminal payoff \(\phi\) in a representation window and let
\(\mathfrak H\) be a compact family of deterministic shifts \(h(t,x)\) for
which the shifted equations in \cref{def:hamiltonian_shift} remain in the same
regularity window.  For each \(h\in\mathfrak H\), let \(u^h\) be the shifted
solution with terminal value \(\phi\).  The calibrated stress envelope is
\[
  \Pi_{\mathfrak H}\phi(t,x)
  =
  \sup_{h\in\mathfrak H} u^h(t,x).
\]
This envelope is a post-derivation stress test for the already constructed
theta-expectation semigroup, not a new primitive market model.
\end{definition}

\begin{proposition}[Robust-pricing stress test]
\label{prop:robust_pricing_stress_test}
In the smooth nondegenerate representation window, each \(h\in\mathfrak H\)
has a payoff-calibrated diffusion and BSDE representation obtained from the
shifted coefficients \(D^h,H^h\).  The envelope
\(\Pi_{\mathfrak H}\phi\) therefore records a family of calibrated drift and
Hamiltonian stresses for the fixed payoff.  When \(D\) is independent of \(p\)
and \(H\) is linear-quadratic, this reduces to the familiar drift-uncertainty
stress test after the Cameron-Martin--Girsanov tilt.  In the nonconvex
theta-expectation case the envelope should not be read as a sublinear
expectation unless the additional sublinearity inequalities are separately
verified.
\end{proposition}

\begin{proof}
For each fixed \(h\), \cref{thm:nonlinear_cm_formula} rewrites the shifted
equation as the unshifted equation plus the deterministic defect
\(\Lambda_h\).  Once the payoff and shifted decoupling field \(u^h\) are
fixed, the construction of
\cref{def:calibrated_linear_generator,thm:smooth_diffusion_representation,thm:fbsde_representation}
applies with \(D,H,u\) replaced by \(D^h,H^h,u^h\).  Compactness of
\(\mathfrak H\) and the representation-window bounds keep the family in one
local regularity class.  In the linear-quadratic case the defect is the
classical drift tilt \(p^\top a h\).  For a nonconvex Hamiltonian, taking a
supremum over calibrated stresses gives a useful pricing diagnostic, but it
does not restore the global subadditivity that Paper 2 proves can fail.
\end{proof}

\begin{proposition}[Post-derivation calibration]
\label{prop:post_derivation_calibration}
The calibrated laws and BSDEs in
\cref{thm:smooth_diffusion_representation,thm:regularized_diffusion_representation,thm:fbsde_representation}
are downstream representations of a fixed payoff.  They do not form a single
payoff-independent martingale problem and cannot be used to prove
\cref{ass:theta_semigroup_input}.
\end{proposition}

\begin{proof}
The law in \cref{thm:smooth_diffusion_representation} depends on \(u\), and
\(u\) is determined by the terminal payoff through the nonlinear HJB semigroup.
Changing the payoff changes \(u\), hence changes \(a^u\), \(b^u\), the
calibrated diffusion, and the BSDE driver.  A payoff-independent law would
contradict \cref{thm:linearity_obstruction} except in accidentally linear
cases.  Therefore the logical order is acyclic: Paper 2 first supplies the
semigroup, then the present paper represents its fixed-payoff derivatives.
\end{proof}

\section{Main Representation Package}

\begin{theorem}[Representation calculus package]
\label{thm:representation_package}
Under \cref{ass:theta_semigroup_input}, every sufficiently regular terminal
payoff determines:
\begin{enumerate}[label=\textup{(\roman*)},wide,labelindent=0pt]
  \item a payoff-calibrated linearized generator $\Lcal^u$;
  \item a carre-du-champ $\Gamma^u$;
  \item a smooth or regularized calibrated diffusion representation;
  \item a decoupled BSDE representation on nondegenerate jet regions;
  \item a semigroup-level Cameron-Martin--Girsanov formula.
  \item a terminal-value sign-convention ledger;
  \item for any compact deterministic stress family, a calibrated
  robust-pricing stress envelope.
\end{enumerate}
All objects above are post-derivation representations of the theta-expectation
semigroup.  None is used to prove the deterministic billiard response theorem
or the HJB homogenization theorem.
\end{theorem}

\begin{proof}
Items (i) and (ii) are \cref{def:calibrated_linear_generator}.  Item (iii) is
\cref{thm:smooth_diffusion_representation,thm:regularized_diffusion_representation}.
Item (iv) is \cref{thm:fbsde_representation}.  Item (v) is
\cref{thm:nonlinear_cm_formula}.  Item (vi) is
\cref{prop:parabolic_orientation_ledger}.  Item (vii) is
\cref{def:calibrated_stress_family,prop:robust_pricing_stress_test}.  The
downstream character follows from the single input assumption
\cref{ass:theta_semigroup_input} and the calibration guardrail in
\cref{prop:post_derivation_calibration}.
\end{proof}

\section{Submission Scope}

The paper's submission version should keep smoothness, nondegeneracy, and
localization assumptions close to each representation theorem.  It should not
import billiard technical lemmas except through Paper 2's HJB theorem.  The
robust-pricing example above should remain a short illustration of calibrated
stress testing, not an additional primitive derivation.

\bibliographystyle{plain}
\bibliography{../../reference}

\end{document}
